Если же в инвариантном подпространстве $L_1$, базисом является $e_{n-k+1}, \dots, e_{n}$, то матрица $A_{\phi}$ имеет вид:
$$
A_{\phi}=\left(
    \begin{array}{c|c}
    A' & 0 \\
    \hline
     A'' & A_{k\times k}
    \end{array}
\right)
$$

Обобщая эти два утверждения, получаем следующий факт:
\begin{theorem}
    Пространство $L$ является прямой суммой инвариантных подпространств $L_1 \oplus L_2 \oplus \dots \oplus L_m$ тогда и только тогда, когда матрица отображения имеет вид 
    $$
        \begin{pmatrix}
        A_1 & 0 & \cdots & 0 \\
        0 & A_2 & \cdots & 0 \\
        \vdots & \vdots & \ddots & \vdots \\
        0 & 0 & \cdots & A_m
        \end{pmatrix}
    $$
где $A_i$ матрица ограничения отображения на $L_i$ (клеточно-диагональная матрица).
\end{theorem}

\section{Собственные векторы}
\subsection{Перестановочные преобразования}
В данном разделе мы иногда будем одинаково обозначать преобразование и его матрицу, смысл обозначения будет зависеть от контекста.
\begin{theorem}
    Пусть $A, B$ — перестановочные линейные преобразования (т.е. композиция не зависит от порядка применения). Тогда ядро и образ одного из преобразований инвариантны относительно другого.
\end{theorem}
\begin{proof}
    \begin{itemize}
        \item Если $x \in \ker A$, то по определению $A(x)=0$, а значит $B(A(x))=0$, но в силу перестановочности $B(A(x))=A(B(x))=0$, следовательно $B(x) \in \ker A$.
        \item Если $x \in \text{Im}(A)$, то по определению $\exists y \hookrightarrow A(y)=x$, а значит $B(x)=B(A(y))=A(B(y)) \in \text{Im}(A)$, следовательно $B(x)\in \text{Im}(A)$.
    \end{itemize}
\end{proof}

Заметим, что в силу перестановочности $A$ и $A^k$, $A$ и любой многочлен $P(A)$ перестановочные. В частности $A$ и $A-\lambda E$, где $\lambda \in \mathbb{R}$, перестановочные; следовательно, по теореме 6.1 $\ker (A-\lambda E)$ инвариантно относительно $A$.

Посмотрим на векторы, принадлежащие $\ker (A-\lambda E)$:
\begin{definition}
    Пусть для некоторого преобразования $A$ и $\lambda \in \mathbb{R}$ ядро $\ker (A-\lambda E)$ ненулевое. Тогда число $\lambda$ называется \textit{собственным значением} преобразования $A$, а $\ker( A-\lambda E)$ \textit{собственным подпространством} относительно преобразования $A$, соответствующее собственному значению $\lambda$.
\end{definition}
Пусть дано преобразование $A$; $\lambda$ — собственное значение; $\ker (A-\lambda E)$ — собственное подпространство. Тогда: 
\begin{enumerate}
    \item Если $\lambda = 0$, то $\ker A$ — собственное подпространство, а ограничение преобразования на $\ker A$ — нулевое преобразование.
    \item Если $\lambda \neq 0$, то из $x \in \ker (A-\lambda E)$ следует $(A-\lambda E)x=0 \Leftrightarrow Ax=\lambda x$, а значит ограничение преобразования $A$ на этом подпространстве — гомотетия с коэффициентом $k=\lambda$.
\end{enumerate}

\begin{definition}
    Вектор $x$ называется \textit{собственным вектором} преобразования $A$, соответствующим собственному значению $\lambda$, если 
    $$
    1) \ x \neq 0, \ \ \ \ \ \  2) \ Ax=\lambda x
    $$
\end{definition}
\textit{Замечание}: из определения следует, что все собственные векторы — это ненулевые векторы собственных подпространств.

С практической точки зрения: если $\ker( A-\lambda E)$, то, для того чтобы найти собственные векторы, нужно решить систему $( A-\lambda E)x=0$.
\begin{theorem}
    Собственные векторы и только они являются базисными в одномерных подпространствах, инвариантных относительно преобразования.
\end{theorem}
\begin{proof}
Дано преобразование $A$.
    \begin{itemize}
        \item Пусть $x$ -- собственный вектор, соответствующий $\lambda$; $L_1$ -- одномерное подпространство с базисом $x$, а $y \in L_1$. Тогда $y=\alpha x$ и $A(y)=\alpha A(x)=\alpha \lambda x \in L_1$, следовательно, $A(y) \in L_1$, а значит $L_1$ - инвариантно относительно  $A$.
        \item Пусть $L_1$ - одномерное инвариантное относительно $A$ подпространство, а $x$ -- базисный вектор. Тогда из $x \in L_1$ следует $A(x)\in L_1$, откуда получаем, что $A(x)=\lambda x$, т.е. $x$ -- собственный вектор, соответсвующий координате $A(x)$ в $L_1$. 
    \end{itemize}
\end{proof}

\textit{Вывод}: в собственном подпространстве любой ненулевой вектор порождает одномерное инвариантное подпространство.

Если в пространстве  существует базис из собственных векторов, то матрица преобразования будет диагональной:
\[
A = \begin{pmatrix}
a_{11} & 0 & 0 & \cdots & 0 \\
0 & a_{22} & 0 & \cdots & 0 \\
0 & 0 & a_{33} & \cdots & 0 \\
\vdots & \vdots & \vdots & \ddots & \vdots \\
0 & 0 & 0 & \cdots & a_{nn}
\end{pmatrix}, \ \ \ \ A(e_i)=a_{ii}e_i
\]
\subsection{Характеристическое уравнение}

Пусть в пространстве $L$ выбран базис $e^{(n)}$; $A$ -- матрица преобразования. Тогда $\ker (A-\lambda E)$ ненулевое $\Leftrightarrow$ $(A-\lambda E)x=0$ имеет нетривиальное решение $\Leftrightarrow$ $\det (A-\lambda E)=0$ (согласно теореме Крамера).
\begin{definition}
    Уравнение $\det (A-\lambda E)=0$ называется характеристическим уравнением, а решения называются характеристическими числами.
\end{definition}
\textit{Замечание}: характеристические числа $\neq$ собственные значения, так как характеристические числа могут быть комплексными, а собственные значения по определению являются вещественными.